% ----------------------------------
%   Arbitrage theory (Chapter 5)
% ----------------------------------

\begin{frame}
\frametitle{Arbitrage theory. Market, portfolios and arbitrage}
\begin{scriptsize}
\begin{itemize}
\item \textbf{Market}: a riskless asset $dS_0 = r S_0\, dt$ and risky assets
\begin{equation*}
    dS_i(t) = \mu_i(S(t), t)\, dt + \sigma_i(S(t), t)^\intercal dW(t), \qquad i = 1, \dots, n .
\end{equation*}
\item \textbf{Portfolio}: an adapted process $\theta(t)$ with value $V(t) = \theta(t)^\intercal S(t)$,
self-financing if $dV(t) = \theta(t)^\intercal dS(t)$.

\item \textbf{PDE method}: replicating $V(S(t), t)$ and imposing self-financing gives
$\theta_i = \partial V / \partial S_i$ and
\begin{equation*}
    \frac{\partial V}{\partial t} + \frac{1}{2} \sum_{i,j} \frac{\partial^2 V}{\partial S_i \partial S_j} \sigma_i^\intercal \sigma_j = 0,
    \qquad V = \sum_{i} \frac{\partial V}{\partial S_i} S_i,
\end{equation*}
which recovers the BS PDE for $n = 1$, $\sigma_1 = \sigma S_1$.

\item \textbf{Arbitrage}: a self-financing admissible portfolio with
\begin{enumerate}
\item $V(0) < 0$ and $\mathbb{P}(V(t) \geq 0) = 1$, or
\item $V(0) = 0$, $\mathbb{P}(V(t) \geq 0) = 1$ and $\mathbb{P}(V(t) > 0) > 0$.
\end{enumerate}
\end{itemize}
\smallskip
Bj\"ork, T. (2009). \emph{Arbitrage Theory in Continuous Time}. Oxford University Press.\\
Glasserman, P. (2003). \emph{Monte Carlo Methods in Financial Engineering}. Springer.
\end{scriptsize}
\end{frame}

\begin{frame}
\frametitle{Arbitrage theory. Martingale method}
\begin{scriptsize}
\begin{itemize}
\item A \textbf{risk-neutral measure} $\Q \sim \mathbb{P}$ is one under which the normalized prices
$Z_i = S_i / S_0$ are martingales: every asset drifts at the risk-free rate.

\item \textbf{First fundamental theorem.}
The market is arbitrage-free if and only if there exists an equivalent (local) martingale measure $\Q$.

\item \textbf{Second fundamental theorem.}
Assuming no arbitrage, the market is complete if and only if $\Q$ is unique.

\item In a complete market the price of a claim with payoff $f(S(T))$ is
\begin{equation*}
    V(t) = \E^{\Q} \left[ \left. e^{ - \int^T_t r(s)\, ds} f(S(T)) \,\right| \F_t \right].
\end{equation*}
\end{itemize}
\medskip
For a call with constant $r$,
\begin{equation*}
    C(K, T) = e^{-rT}\, \E^{\Q}\!\left[(S_T - K)^+\right],
\end{equation*}
which is the starting point of the Dupire equation.
\end{scriptsize}
\end{frame}
